\documentclass[../../../include/open-logic-section]{subfiles}

\begin{document}

\olfileid[id]{sth}{z}{milestone}
\olsection{$\Z^-$: Sebuah Tonggak Penting}

Kita akan kembali membahas \stagesinf{} dalam bagian berikutnya. Namun,
dengan Aksioma Ketakhinggaan, kita telah mencapai tonggak penting. Sekarang
kita mempunyai semua aksioma yang diperlukan bagi teori~$\Zminus$. Secara
terperinci:

\begin{defn}
Teori $\Zminus$ mempunyai aksioma-aksioma berikut: Ekstensionalitas, Gabungan,
Pasangan, Himpunan Kuasa, Ketakhinggaan, dan semua instans skema Pemisahan.
\end{defn}

Nama itu merupakan singkatan bagi teori himpunan \emph{Zermelo}
(\emph{minus} sesuatu yang akan kita jumpai kemudian). Zermelo layak menerima
kehormatan ini karena ia pada dasarnya merumuskan teori tersebut pada
\citeyear{Zermelo1908Untersuchungen}.\footnote{Untuk komentar menarik mengenai
sejarah dan rincian teknisnya, lihat \citet[Lampiran
A]{Potter2004}.}

Teori ini cukup kuat untuk memungkinkan kita mengerjakan sangat banyak
matematika. Secara khusus, Anda \emph{sebaiknya} melihat kembali seluruh
\olref[sfr][][]{part} dan meyakinkan diri bahwa segala sesuatu yang kita
kerjakan secara naif dapat dikerjakan secara lebih formal di dalam~$\Zminus$.
(Setelah melakukannya cukup lama, Anda mungkin ingin melompat ke depan dan
membaca \olref[sth][z][arbintersections]{sec}.) Jadi, mulai sekarang dan tanpa
komentar lebih lanjut, kita menganggap diri kita bekerja sekurang-kurangnya
di dalam~$\Zminus$.

\end{document}
